\chapter{Integration by Parts}

In the last chapter we learned to undo the chain rule. Substitution let us
recognize when an integrand was secretly born from a composition $f(g(x))$ and to
peel the composition back apart. That handles an enormous number of integrals --
but it leaves a whole continent untouched. What do you do with
\[
\int x\,e^{x}\,\dd x, \qquad \int x\sin x\,\dd x, \qquad \int \ln x\,\dd x, \qquad
\int e^{x}\cos x\,\dd x?
\]
None of these is a chain rule in disguise. Try substitution on
$\int x e^{x}\,\dd x$ and you will flail: setting $u=x$ does nothing, setting
$u=e^{x}$ leaves a stray $x$ you cannot get rid of. The trouble is that these
integrands are \emph{products} of two unrelated things -- a polynomial times an
exponential, a polynomial times a sine -- and substitution was built to undo
compositions, not products.

There is exactly one differentiation rule that creates products: the
\emph{product rule}. So if we want to integrate products, the move is to
run the product rule backwards, exactly as substitution was the chain rule run
backwards. The technique that comes out is called \emph{integration by parts},
and it is the second of the two great workhorses of integration. It will not
collapse an integral to nothing the way a lucky substitution sometimes does.
Instead it does something subtler and, once you see it, rather beautiful: it
\emph{trades} the integral you have for a different one, in the hope that the new
one is easier. Most of the art is in arranging the trade so that you come out
ahead.

\chapterline

\section{The Product Rule, Run Backwards}
\label{ch23:product-backwards}

Everything in this chapter is squeezed out of one line you already know. If $u$
and $v$ are differentiable functions of $x$, the product rule says
\[
\ddx{u\,v} = u\,\dv{v}{x} + v\,\dv{u}{x}.
\]
Read it as a statement about how the area of a little rectangle $u\cdot v$
changes: when both sides are growing, the rectangle grows on two fronts at once,
one strip of width $\dd u$ down the side and one strip of height $\dd v$ along
the top. The product rule just adds those two strips. (We drew that picture back
when we first met the product rule; hold it in mind, because it is about to pay
off.)

Now integrate both sides with respect to $x$ over the whole real line of the
problem. Integration undoes differentiation, so the left side integrates to
$u\,v$ itself:
\[
u\,v = \int u\,\dv{v}{x}\,\dd x + \int v\,\dv{u}{x}\,\dd x.
\]
That is already the whole theorem -- it just needs to be rearranged. Solve for
the first integral on the right:
\[
\int u\,\dv{v}{x}\,\dd x = u\,v - \int v\,\dv{u}{x}\,\dd x.
\]
Finally, abbreviate $\dd v = \dv{v}{x}\,\dd x$ and $\dd u = \dv{u}{x}\,\dd x$ --
the differentials of Chapter~14, used here as plain bookkeeping for ``the slope
times the step.'' We arrive at the formula every calculus student eventually
recites in their sleep:

\begin{focusbox}
\textbf{Integration by parts.} If $u$ and $v$ are differentiable, then
\[
\boxed{\;\int u\,\dd v \;=\; u\,v \;-\; \int v\,\dd u\;}
\]
and, between limits $a$ and $b$,
\[
\int_{a}^{b} u\,\dd v \;=\; \eval{u\,v}{a}^{b} \;-\; \int_{a}^{b} v\,\dd u.
\]
\end{focusbox}

That is it. Two symbols, $u$ and $v$, and the rule says: integrating $u\,\dd v$
costs you the boundary term $uv$ minus the price of integrating $v\,\dd u$
instead. You have not solved the integral -- you have \emph{swapped} it for
another. The bet is that $\int v\,\dd u$ is friendlier than $\int u\,\dd v$.

\begin{whybox}
Substitution is the chain rule read right to left; integration by parts is the
product rule read right to left. That duality organizes the
whole structure of integration technique. There are only a handful of ways to
\emph{build} a function by differentiation (sums, products, compositions), so
there are only a handful of ways to \emph{take one apart} by integration. Sums
split term-by-term. Compositions are unwound by substitution. Products are unwound
by parts. When you face an unfamiliar integral, the first question is always
``which differentiation rule could have produced an integrand shaped like this?''
-- and the answer tells you which tool to reach for.
\end{whybox}

\begin{happeningbox}
Look again at $\int u\,\dd v = uv - \int v\,\dd u$ and ask what is physically
being moved around. The combination $u\,\dd v$ is a thin horizontal slab of
``stuff'': a width $u$ times an increment $\dd v$. The formula says the total of
all those strips equals the finished rectangle $uv$ \emph{minus} the total of the
complementary strips $v\,\dd u$. Integration by parts is thus a
statement about cutting a rectangle two ways. The product $uv$ is a region in the
$u$--$v$ plane; you can sweep it out in horizontal slabs ($u\,\dd v$) or vertical
slabs ($v\,\dd u$), and the two sweeps must add up to the same rectangle. We will
make this picture exact in the next section, and it explains \emph{why} the
boundary term $uv$ is there.
\end{happeningbox}

\section{The Picture: Slicing a Rectangle Two Ways}
\label{ch23:picture}

The boundary term $uv$ can feel like it falls out of the algebra by accident. It
does not. Here is a proof of integration by parts that uses no symbols at all
until the very end -- just the area of a rectangle.

Imagine a curve in the $u$--$v$ plane traced out as $x$ runs from $a$ to $b$. At
$x=a$ the point is $(u_a, v_a)$; at $x=b$ it is $(u_b, v_b)$. As $x$ increases,
both $u$ and $v$ change, and the point sweeps a path. Now look at
Figure~\ref{ch23:fig-rectangle}. The big rectangle reaching out to $(u_b,v_b)$
has area $u_b v_b$. The small rectangle reaching to $(u_a,v_a)$ has area
$u_a v_a$. The L-shaped region between them -- the big rectangle minus the small
one -- has area $u_b v_b - u_a v_a$, which is exactly the boundary term
$\eval{uv}{a}^{b}$.

\begin{figure}[H]
\centering
\begin{tikzpicture}[scale=1.1]
% Axes
\draw[->, thick] (0,0) -- (6.8,0) node[right, font=\small] {$u$};
\draw[->, thick] (0,0) -- (0,5.4) node[above, font=\small] {$v$};

% Corner points
\coordinate (O) at (0,0);
\coordinate (A) at (1.6,1.2);
\coordinate (B) at (5.4,4.4);

% Shaded inner rectangle (u_a * v_a)
\fill[gray!15] (0,0) rectangle (1.6,1.2);
\draw[gray!60, thin] (0,0) rectangle (1.6,1.2);
\node[font=\footnotesize, gray!70!black] at (0.8,0.6) {$u_a v_a$};

% Smooth curve
\def\curvepath{(1.6,1.2) .. controls (2.6,2.2) and (3.8,3.2) .. (5.4,4.4)}

% Region 1: horizontal slabs \int u dv (left/above curve)
\fill[sourceblue!20] (0,1.2) -- (1.6,1.2) .. controls (2.6,2.2) and (3.8,3.2) .. (5.4,4.4) -- (0,4.4) -- cycle;

% Region 2: vertical slabs \int v du (below/right curve)
\fill[quotientorange!20] (1.6,0) -- (1.6,1.2) .. controls (2.6,2.2) and (3.8,3.2) .. (5.4,4.4) -- (5.4,0) -- cycle;

% Draw curve
\draw[curveblue, line width=1.5pt] \curvepath;

% Outer bounding box
\draw[dashed, gray!70] (5.4,0) -- (5.4,4.4) -- (0,4.4);
\draw[dashed, gray!70] (1.6,0) -- (1.6,1.2) -- (0,1.2);

% Sample horizontal strip in Region 1
\fill[sourceblue!60] (0,2.8) rectangle (3.3,2.95);
\draw[<->, sourceblue!90!black, thin] (0,2.6) -- (3.3,2.6) node[midway, below, font=\tiny] {width $u$};
\draw[<->, sourceblue!90!black, thin] (-0.2,2.8) -- (-0.2,2.95) node[midway, left, font=\tiny] {$\dd v$};

% Sample vertical strip in Region 2
\fill[quotientorange!70] (4.0,0) rectangle (4.15,3.35);
\draw[<->, quotientorange!90!black, thin] (4.3,0) -- (4.3,3.35) node[midway, right, font=\tiny] {height $v$};
\draw[<->, quotientorange!90!black, thin] (4.0,-0.2) -- (4.15,-0.2) node[midway, below, font=\tiny] {$\dd u$};

% Endpoints
\fill[curveblue] (A) circle (2pt) node[above left, font=\footnotesize] {$(u_a,v_a)$};
\fill[curveblue] (B) circle (2pt) node[above right, font=\footnotesize] {$(u_b,v_b)$};

% Axis tick marks
\node[below, font=\footnotesize] at (1.6,0) {$u_a$};
\node[below, font=\footnotesize] at (5.4,0) {$u_b$};
\node[left, font=\footnotesize] at (0,1.2) {$v_a$};
\node[left, font=\footnotesize] at (0,4.4) {$v_b$};

% Region labels (placed in clear whitespace)
\node[font=\small\bfseries, sourceblue!80!black] at (1.5,3.8) {$\displaystyle\int_{v_a}^{v_b} u\,\dd v$};
\node[font=\small\bfseries, quotientorange!80!black] at (4.2,1.5) {$\displaystyle\int_{u_a}^{u_b} v\,\dd u$};

\end{tikzpicture}
\caption{Integration by parts as bounding box geometry. The total rectangle out to $(u_b,v_b)$ has area $u_b v_b$. Removing the bottom-left rectangle $u_a v_a$ leaves an L-shaped region of area $u_b v_b - u_a v_a = \eval{uv}{a}^b$. The path $(u(x),v(x))$ cuts this L-region into two complementary sweeps: horizontal strips $u\,\dd v$ (blue) summing to $\int u\,\dd v$, and vertical strips $v\,\dd u$ (orange) summing to $\int v\,\dd u$. Together they tile the region: $\int u\,\dd v + \int v\,\dd u = \eval{uv}{a}^b$.}
\label{ch23:fig-rectangle}
\end{figure}

Now slice that L-shaped region two different ways. A \emph{vertical} strip has a
height $v$ and is built from a change in $u$; its area is $v\,\dd u$, and summing
them gives $\int v\,\dd u$, the orange region below the curve. A \emph{horizontal}
strip has a width $u$ and is built from a change in $v$; its area is $u\,\dd v$,
and summing gives $\int u\,\dd v$, the blue region to the left of the curve. The
path of $(u,v)$ is the boundary between the two regions, and together they tile
the entire L. Therefore
\[
\int_a^b u\,\dd v \;+\; \int_a^b v\,\dd u \;=\; u_b v_b - u_a v_a
= \eval{uv}{a}^{b}.
\]
Move the second integral to the right and you have integration by parts, born
purely from cutting a region two ways.

\begin{intuitionbox}
This is why the boundary term is $uv$ and not something else: $uv$ is literally
the \emph{rectangle} whose two complementary sweeps are the two integrals.
Integration by parts is conservation of area. Whatever you ``give up'' by not
integrating $u\,\dd v$ directly, you get back as the rectangle minus the other
sweep. Nothing is created or destroyed; the area is just being cut into different
strips.
\end{intuitionbox}

This geometric proof needs $u$ and $v$ to be monotonic to draw cleanly, but the
algebraic derivation in Section~\ref{ch23:product-backwards} has no such
restriction -- it works for any differentiable $u,v$. Keep the picture for
intuition and the algebra for rigor.

\section{Choosing \texorpdfstring{$u$}{u} and \texorpdfstring{$dv$}{dv}}
\label{ch23:choosing}

The formula has a built-in decision. Faced with an integral, you must split the
integrand into two pieces: one piece you will call $u$ (and \emph{differentiate}
to get $\dd u$), and one piece you will call $\dd v$ (and \emph{integrate} to get
$v$). Every product can be split two ways, and usually only one split helps. So
how do you choose?

Two principles guide every good choice.

\begin{enumerate}
\item \textbf{$u$ should get simpler when you differentiate it.} The whole point
is that $\int v\,\dd u$ is easier than $\int u\,\dd v$. Differentiating $u$ is
your chance to make progress, so pick a $u$ whose derivative is tamer than
itself. A polynomial like $x^3$ is ideal: each derivative knocks the degree down,
and after enough rounds it vanishes. The logarithm $\ln x$ is even better -- it
has no elementary antiderivative of its own, but its \emph{derivative} $1/x$ is
beautifully simple, so making $\ln x$ into $u$ converts a hard factor into an
easy one.

\item \textbf{$\dd v$ must be something you can actually integrate.} You need
$v=\int \dd v$ to be a function you know. There is no point assigning $\dd v$ to a
factor you cannot antidifferentiate. So $e^x\,\dd x$, $\sin x\,\dd x$,
$\cos x\,\dd x$, and powers $x^n\,\dd x$ all make fine $\dd v$'s; something like
$\ln x\,\dd x$ does not (you would be stuck integrating it -- which is the very
problem you started with).
\end{enumerate}

These two principles together produce a famous mnemonic.

\begin{focusbox}
\textbf{LIATE} -- a priority order for choosing $u$. When a product mixes two of
these types, let $u$ be whichever appears \emph{earlier} in the list, and let the
rest be $\dd v$:
\begin{center}
\begin{tabular}{cl}
\textbf{L} & \textbf{L}ogarithmic ($\ln x$, $\log x$) \\
\textbf{I} & \textbf{I}nverse trig ($\arctan x$, $\arcsin x$) \\
\textbf{A} & \textbf{A}lgebraic (powers $x^n$, polynomials) \\
\textbf{T} & \textbf{T}rigonometric ($\sin x$, $\cos x$) \\
\textbf{E} & \textbf{E}xponential ($e^x$, $a^x$) \\
\end{tabular}
\end{center}
Earlier in the list $=$ choose as $u$.
\end{focusbox}

Why this order? Read the list top to bottom and you are reading functions in
order of ``how much simpler does differentiating make them, and how easy are they
to integrate.'' Logs and inverse trig functions are at the top because they are
miserable to integrate but pleasant to differentiate, so you always want them as
$u$. Exponentials are at the bottom because they are trivial to integrate and
never get simpler when differentiated, so you always want them as $\dd v$.
Polynomials sit in the middle: they are happy either way, but differentiating
lowers their degree, so when paired with a trig or exponential they go on the $u$
side.

\begin{trapbox}
LIATE is a \emph{heuristic, not a law}. It encodes good defaults, but it cannot
think for you, and there are integrals where it points the wrong way or where
\emph{either} choice works. The genuine criterion is always the two principles
above: did $\int v\,\dd u$ come out simpler than what you started with? If a
choice leads to a worse integral, you have learned something -- swap your choice
and try the other split. Never defend a bad choice by appeal to a mnemonic.
\end{trapbox}

\begin{historybox}
The acronym LIATE is a twentieth-century pedagogical heuristic -- introduced in a 1983
paper by Herbert Kasube to assist calculus students. The foundational \emph{method}
of integration by parts has a far grander lineage: it was systematically introduced by
the English mathematician \emph{Brook Taylor} in his 1715 treatise \emph{Methodus
Incrementorum Directa et Inversa} (the same work that gave birth to Taylor series).
Taylor recognized that the Product Rule of differentiation:
\[
  \dd(uv) = u\,\dd v + v\,\dd u
\]
could be integrated and rearranged to shift a derivative off an intractable factor and onto
a compliant one.

In modern mathematics and physics, integration by parts is not just an integration trick;
it is a structural bridge across disciplines. In nineteenth-century analysis, Pierre-Simon
\emph{Laplace} and Adrien-Marie \emph{Legendre} recognized parts as the fundamental engine of
integral transforms. In the Laplace transform $\mathcal{L}\{f'\}(s) = \int_0^\infty e^{-st} f'(t)\,\dd t$,
a single integration by parts yields $s\mathcal{L}\{f\}(s) - f(0)$, converting intractable
differential equations directly into high-school algebra. In Hamiltonian mechanics and thermodynamics,
the Legendre transform converts the Lagrangian $L(q, \dot{q})$ into the Hamiltonian $H(q, p) = p\dot{q} - L$
by swapping coordinates across the complementary bounding box $uv = \int u\,\dd v + \int v\,\dd u$.
And in quantum mechanics and modern partial differential equations, integration by parts is how
mathematicians define \emph{weak derivatives} and Hermitian operators, moving derivatives
off un-differentiable physical wavefunctions onto smooth test functions.
Leibniz's notation $\int u\,\dd v = uv - \int v\,\dd u$ is what made this universal principle
effortless to write and remember.
\end{historybox}

\section{First Worked Examples}
\label{ch23:examples}

The ritual is always the same four steps: choose $u$ and $\dd v$; differentiate
to get $\dd u$ and integrate to get $v$; assemble $uv - \int v\,\dd u$; finish
the leftover integral.

\begin{example}[The archetype: $\int x e^{x}\,\dd x$]
\label{ch23:ex-xex}
Here the two factors are algebraic ($x$) and exponential ($e^x$). LIATE says: A
beats E, so $u=x$ and $\dd v = e^x\,\dd x$. Then
\[
u = x \;\Rightarrow\; \dd u = \dd x, \qquad
\dd v = e^x\,\dd x \;\Rightarrow\; v = e^x.
\]
Assemble:
\[
\int x e^{x}\,\dd x = uv - \int v\,\dd u = x e^{x} - \int e^{x}\,\dd x
= x e^{x} - e^{x} + C = (x-1)e^{x} + C.
\]
Notice what happened: the troublesome factor $x$ got differentiated down to $1$,
turning the leftover integral $\int e^x\,\dd x$ into something we can do in our
sleep. \emph{That} is the payoff of choosing $u=x$.
\end{example}

\begin{labbox}
Check Example~\ref{ch23:ex-xex} by differentiating the answer. By the product
rule, $\ddx{(x-1)e^x} = 1\cdot e^x + (x-1)e^x = e^x + xe^x - e^x = xe^x$. It
matches the integrand. \emph{Always} verify a by-parts answer this way at first --
differentiation is fast and it catches sign errors instantly. Integration by
parts is sign-error-prone precisely because of the minus sign in the formula.
\end{labbox}

\begin{trapbox}
What if you had chosen the other way, $u=e^x$ and $\dd v = x\,\dd x$? Then
$\dd u = e^x\,\dd x$ and $v = x^2/2$, and the formula gives
\[
\int x e^x\,\dd x = \frac{x^2}{2}e^x - \int \frac{x^2}{2}e^x\,\dd x.
\]
The new integral has $x^2$ where the old one had $x$ -- it got \emph{worse}. The
polynomial went up in degree because we integrated it instead of differentiating
it. This is the classic wrong turn, and it is exactly what LIATE warns you off:
never integrate the polynomial when there is an exponential or trig factor
waiting to absorb the $\dd v$ role.
\end{trapbox}

\begin{example}[The clever trick: $\int \ln x\,\dd x$]
\label{ch23:ex-lnx}
This one looks impossible -- there is no second factor to play with. The trick is
to \emph{invent} one by writing $\ln x = (\ln x)\cdot 1$. Now it is a product, and
LIATE is decisive: L beats everything, so $u = \ln x$ and $\dd v = 1\,\dd x$.
Then
\[
u = \ln x \;\Rightarrow\; \dd u = \frac{1}{x}\,\dd x, \qquad
\dd v = \dd x \;\Rightarrow\; v = x.
\]
Assemble:
\[
\int \ln x\,\dd x = (\ln x)(x) - \int x\cdot\frac{1}{x}\,\dd x
= x\ln x - \int 1\,\dd x = x\ln x - x + C.
\]
The miracle is in the cancellation $x\cdot\frac{1}{x}=1$: differentiating $\ln x$
produced a $1/x$ that annihilated the $x$ from $v$, leaving the trivial integral
$\int 1\,\dd x$. This is the standard way to integrate \emph{any} function that is
easy to differentiate but has no obvious antiderivative -- log, $\arctan$,
$\arcsin$ all yield to ``$u=$ the function, $\dd v = \dd x$.''
\end{example}

\begin{example}[Inverse trig: $\int \arctan x\,\dd x$]
\label{ch23:ex-arctan}
Same trick. Write $\arctan x = (\arctan x)\cdot 1$, take $u=\arctan x$ and
$\dd v = \dd x$, so $\dd u = \dfrac{1}{1+x^2}\,\dd x$ and $v=x$. Then
\[
\int \arctan x\,\dd x = x\arctan x - \int \frac{x}{1+x^2}\,\dd x.
\]
The leftover integral now \emph{is} a substitution problem -- set $w = 1+x^2$,
$\dd w = 2x\,\dd x$, so $\int \frac{x}{1+x^2}\,\dd x = \frac12\ln(1+x^2)$. Hence
\[
\int \arctan x\,\dd x = x\arctan x - \tfrac12\ln(1+x^2) + C.
\]
Notice how the two big techniques chain together: parts cracked the problem open,
and substitution mopped up the remainder. Real integration is almost always a
relay race between methods.
\end{example}

\begin{example}[A polynomial that needs two rounds: $\int x^2 \sin x\,\dd x$]
\label{ch23:ex-x2sin}
Here $u=x^2$ (algebraic, beats trig) and $\dd v = \sin x\,\dd x$, so
$\dd u = 2x\,\dd x$ and $v = -\cos x$. One round:
\[
\int x^2\sin x\,\dd x = -x^2\cos x - \int(-\cos x)(2x\,\dd x)
= -x^2\cos x + 2\int x\cos x\,\dd x.
\]
We have not finished, but we have made progress: the power dropped from $x^2$ to
$x$. Attack $\int x\cos x\,\dd x$ by parts again, with $u=x$,
$\dd v = \cos x\,\dd x$, $\dd u = \dd x$, $v = \sin x$:
\[
\int x\cos x\,\dd x = x\sin x - \int \sin x\,\dd x = x\sin x + \cos x.
\]
Substitute back:
\[
\int x^2\sin x\,\dd x = -x^2\cos x + 2\big(x\sin x + \cos x\big) + C
= -x^2\cos x + 2x\sin x + 2\cos x + C.
\]
Each application of parts peeled off one degree of the polynomial. A degree-$n$
polynomial times $\sin x$ or $e^x$ surrenders after exactly $n$ rounds. When you
see that coming, there is a shortcut -- tabular integration -- which we meet in
Section~\ref{ch23:tabular}.
\end{example}

\begin{figure}[H]
\centering
\begin{tikzpicture}
\begin{axis}[
  width=11cm, height=6cm,
  xmin=-6.3, xmax=6.3, ymin=-25, ymax=36,
  xtick={-6.28,-3.14,0,3.14,6.28},
  xticklabels={$-2\pi$,$-\pi$,$0$,$\pi$,$2\pi$},
  ytick={-20,-10,0,10,20},
  axis lines=middle,
  tick label style={font=\footnotesize},
  legend style={at={(0.01,0.99)}, anchor=north west, font=\scriptsize, draw=none, fill=none},
]
\addplot[curveblue, thick, samples=180, domain=-6.28:6.28] {x^2*sin(deg(x))};
\addlegendentry{$x^2\sin x$ (integrand)}
\addplot[tangentred, thick, samples=180, domain=-6.28:6.28] {-x^2*cos(deg(x)) + 2*x*sin(deg(x)) + 2*cos(deg(x))};
\addlegendentry{$-x^2\cos x + 2x\sin x + 2\cos x$ (antiderivative)}
\end{axis}
\end{tikzpicture}
\caption{The integrand $x^2\sin x$ (blue) and the antiderivative we computed in
Example~\ref{ch23:ex-x2sin} (red). Wherever the blue curve is positive the red
curve climbs, and where blue is negative red falls -- the antiderivative
accumulates signed area under the integrand, exactly as the Fundamental Theorem
promises.}
\label{ch23:fig-x2sin}
\end{figure}

\begin{example}[Definite integral: $\int_1^{e} \ln x\,\dd x$]
\label{ch23:ex-deflnx}
We already found the antiderivative $x\ln x - x$ in Example~\ref{ch23:ex-lnx}, so
\[
\int_1^{e} \ln x\,\dd x = \eval{x\ln x - x}{1}^{e}
= (e\ln e - e) - (1\cdot\ln 1 - 1)
= (e - e) - (0 - 1) = 1.
\]
Geometrically: the area under the curve $y=\ln x$ from $x=1$ to $x=e$ is exactly
$1$ square unit (Figure~\ref{ch23:fig-lnarea}). That clean answer is no
accident -- $\ln x$ and $e$ are bound together, and the region under one
natural-log arch from $1$ to $e$ has unit area.
\end{example}

\begin{figure}[H]
\centering
\begin{tikzpicture}
\begin{axis}[
  width=10cm, height=5.6cm,
  xmin=0, xmax=3.4, ymin=-0.2, ymax=1.3,
  xtick={1,2,2.718},
  xticklabels={$1$,$2$,$e$},
  ytick={0,0.5,1},
  axis lines=middle,
  tick label style={font=\footnotesize},
  legend pos=north west,
  legend style={font=\footnotesize, draw=none, fill=none},
]
\addplot[curveblue, thick, samples=120, domain=0.37:3.4] {ln(x)};
\addlegendentry{$y=\ln x$}
\addplot[draw=none, name path=curveA, samples=120, domain=1:2.718] {ln(x)};
\addplot[draw=none, name path=axisA, samples=2, domain=1:2.718] {0};
\addplot[areafill, opacity=0.5] fill between[of=curveA and axisA];
\node at (2.05,0.33) {\footnotesize area $=1$};
\end{axis}
\end{tikzpicture}
\caption{The shaded region under $y=\ln x$ from $x=1$ to $x=e$ has area exactly
$1$, the value of $\int_1^e \ln x\,\dd x$ computed in
Example~\ref{ch23:ex-deflnx}. Integration by parts handed us the antiderivative
$x\ln x - x$; the Fundamental Theorem turned it into this clean unit area.}
\label{ch23:fig-lnarea}
\end{figure}

\begin{walkbox}
Here is a question to chew on while walking. We integrated $\ln x$ by parts and
got $x\ln x - x$. But integration by parts came from the product rule, and the
product rule came from differentiation. Yet $\ln x$ is not a product of anything
obvious. How did a rule about products help with a single function? (Hint: the
hidden product is $\ln x \cdot 1$, and the deeper truth is that ``$1$'' is the
derivative of $x$, so we are secretly using $v=x$ as a lever. Where else might
multiplying by a cleverly-chosen $1$ unlock a problem?)
\end{walkbox}

\section{Reduction Formulas: Parts as a Machine}
\label{ch23:reduction}

When the integrand carries a parameter -- a power $n$, say -- a single
application of parts can express the integral with parameter $n$ in terms of the
\emph{same} integral with a smaller parameter. That recursive relationship is
called a \emph{reduction formula}. It is integration by parts at its most
powerful: instead of solving one integral, you solve an entire infinite family in
one stroke.

\begin{example}[The reduction formula for $\int x^{n} e^{x}\,\dd x$]
\label{ch23:ex-redxn}
Let
\[
I_n = \int x^n e^x\,\dd x.
\]
Take $u = x^n$ and $\dd v = e^x\,\dd x$, so $\dd u = n x^{n-1}\,\dd x$ and
$v = e^x$. Parts gives
\[
I_n = x^n e^x - \int e^x \cdot n x^{n-1}\,\dd x
= x^n e^x - n\int x^{n-1}e^x\,\dd x = x^n e^x - n\,I_{n-1}.
\]
So
\[
\boxed{\,I_n = x^n e^x - n\,I_{n-1}.}
\]
Each use of parts lowers the exponent by one. Starting from
$I_0 = \int e^x\,\dd x = e^x$, we can climb the ladder to any $n$. For instance,
\[
I_1 = x e^x - 1\cdot I_0 = xe^x - e^x,
\]
\[
I_2 = x^2 e^x - 2 I_1 = x^2 e^x - 2(xe^x - e^x) = x^2 e^x - 2x e^x + 2e^x,
\]
\[
I_3 = x^3 e^x - 3 I_2 = x^3 e^x - 3x^2 e^x + 6x e^x - 6 e^x,
\]
and the pattern (with its alternating falling-factorial coefficients) is now
unmistakable. One reduction formula generated the whole family.
\end{example}

\begin{whybox}
Why do reduction formulas exist at all? Because parts \emph{trades a hard factor
for a smaller version of itself}. When you differentiate $x^n$ you get
$n x^{n-1}$ -- structurally the same kind of object, one notch simpler. The
integral on the right is therefore a clone of the integral on the left with the
parameter decremented. Anytime differentiating one factor reproduces a scaled,
simpler copy of that factor -- powers, $\sin^n$, $\cos^n$, $(\ln x)^n$ -- parts
will hand you a recursion. The technique is not just for one-off integrals; it is
a way of taming infinite families by induction on a parameter.
\end{whybox}

\begin{example}[A trig reduction: $\int \sin^{n}x\,\dd x$]
\label{ch23:ex-redsin}
This is the classic and it is worth seeing in full, because the cleverness is
real. Let
\[
J_n = \int \sin^n x\,\dd x.
\]
Split off one factor of sine: write $\sin^n x = \sin^{n-1}x\cdot \sin x$, and take
\[
u = \sin^{n-1}x, \qquad \dd v = \sin x\,\dd x,
\]
so that $\dd u = (n-1)\sin^{n-2}x\,\cos x\,\dd x$ and $v = -\cos x$. Parts:
\[
J_n = -\sin^{n-1}x\,\cos x + (n-1)\int \sin^{n-2}x\,\cos^2 x\,\dd x.
\]
Now use the Pythagorean identity $\cos^2 x = 1 - \sin^2 x$ to turn the leftover
integral back into sines:
\[
\int \sin^{n-2}x\,\cos^2 x\,\dd x
= \int \sin^{n-2}x\,\dd x - \int \sin^{n}x\,\dd x = J_{n-2} - J_n.
\]
Substitute:
\[
J_n = -\sin^{n-1}x\,\cos x + (n-1)\big(J_{n-2} - J_n\big).
\]
The integral $J_n$ has \emph{reappeared on the right} -- this is the recurring
integral trick we will name in the next section. Collect the $J_n$ terms:
\[
J_n + (n-1)J_n = -\sin^{n-1}x\,\cos x + (n-1)J_{n-2},
\]
\[
n\,J_n = -\sin^{n-1}x\,\cos x + (n-1)J_{n-2},
\]
and finally
\[
\boxed{\,J_n = -\frac{1}{n}\sin^{n-1}x\,\cos x + \frac{n-1}{n}J_{n-2}.}
\]
This drops the power by \emph{two} each time -- which is exactly right, because
the base cases come in two flavors: $J_0 = \int \dd x = x$ and
$J_1 = \int \sin x\,\dd x = -\cos x$. Every even power reduces down to $J_0$,
every odd power down to $J_1$.
\end{example}

\begin{connectionbox}
The reduction formula for $\int_0^{\pi/2}\sin^n x\,\dd x$ leads to one of the most
beautiful formulas in mathematics. Evaluated between $0$ and $\pi/2$ the boundary
term $\sin^{n-1}x\cos x$ vanishes at both ends, so the recursion becomes purely
numerical: $W_n = \frac{n-1}{n}W_{n-2}$, where $W_n = \int_0^{\pi/2}\sin^n x\,\dd
x$ are the \emph{Wallis integrals}. Chasing the ratios $W_n$ as $n\to\infty$
gives John Wallis's 1655 product
\[
\frac{\pi}{2} = \frac{2}{1}\cdot\frac{2}{3}\cdot\frac{4}{3}\cdot\frac{4}{5}
\cdot\frac{6}{5}\cdot\frac{6}{7}\cdots,
\]
a product of integers that converges to half of $\pi$. That such a thing falls
out of a plain reduction formula is one of the surprises that keeps the subject
interesting.
\end{connectionbox}

\section{The Recurring-Integral Trick}
\label{ch23:recurring}

Sometimes you apply parts, then apply it again, and the integral you started with
\emph{comes back around} -- the very same integral reappears on the right-hand
side. A beginner panics (``I'm going in circles!''). An expert grins, because a
repeated integral is something you can \emph{solve for algebraically}, exactly the
way you'd solve $x = 5 - 2x$ for $x$. We just saw a glimpse of this in the
$\sin^n$ reduction. The cleanest example is the integral that mixes an
exponential with a trig function.

\begin{example}[Solving for the integral: $\int e^{x}\cos x\,\dd x$]
\label{ch23:ex-excos}
Call it $K = \int e^x\cos x\,\dd x$. Neither factor gets simpler when
differentiated -- $e^x$ stays $e^x$, and $\cos x$ cycles to $-\sin x$ to
$-\cos x$ -- so LIATE has no traction. It does not matter which factor we pick for
$u$; let us take $u = e^x$, $\dd v = \cos x\,\dd x$ consistently. Then
$\dd u = e^x\,\dd x$, $v = \sin x$:
\[
K = e^x\sin x - \int e^x\sin x\,\dd x.
\]
The new integral is no simpler, but it is the same \emph{type}. Apply parts again
to $\int e^x\sin x\,\dd x$, keeping the same choice $u = e^x$,
$\dd v = \sin x\,\dd x$, so $v = -\cos x$:
\[
\int e^x\sin x\,\dd x = -e^x\cos x + \int e^x\cos x\,\dd x = -e^x\cos x + K.
\]
Substitute this back into the expression for $K$:
\[
K = e^x\sin x - \big(-e^x\cos x + K\big) = e^x\sin x + e^x\cos x - K.
\]
There it is -- $K$ on both sides. Add $K$ to both sides:
\[
2K = e^x\sin x + e^x\cos x,
\]
\[
\boxed{\,K = \int e^x\cos x\,\dd x = \frac{e^x(\sin x + \cos x)}{2} + C.}
\]
We never ``finished'' an integral in the usual sense -- we trapped the unknown
integral in an equation and solved for it.
\end{example}

\begin{figure}[H]
\centering
\begin{tikzpicture}
\begin{axis}[
  width=11cm, height=6cm,
  xmin=-3.3, xmax=3.3, ymin=-12, ymax=18,
  xtick={-3.14,-1.57,0,1.57,3.14},
  xticklabels={$-\pi$,$-\tfrac{\pi}{2}$,$0$,$\tfrac{\pi}{2}$,$\pi$},
  ytick={-10,0,10},
  axis lines=middle,
  tick label style={font=\footnotesize},
  legend pos=north west,
  legend style={font=\footnotesize, draw=none, fill=none},
]
\addplot[curveblue, thick, samples=200, domain=-3.14:3.14] {exp(x)*cos(deg(x))};
\addlegendentry{$e^x\cos x$ (integrand)}
\addplot[tangentred, thick, samples=200, domain=-3.14:3.14] {exp(x)*(sin(deg(x))+cos(deg(x)))/2};
\addlegendentry{$\tfrac12 e^x(\sin x+\cos x)$ (antiderivative)}
\end{axis}
\end{tikzpicture}
\caption{The integrand $e^x\cos x$ (blue) and the antiderivative
$\tfrac12 e^x(\sin x+\cos x)$ (red) we trapped algebraically in
Example~\ref{ch23:ex-excos}. The exponential factor makes both curves swing
ever wider, but the antiderivative still rises exactly where the integrand is
positive and falls where it is negative -- the recurring-integral trick produced
a genuine antiderivative, not a sleight of hand.}
\label{ch23:fig-excos}
\end{figure}

\begin{trapbox}
The fatal mistake in the recurring trick is to \emph{switch your choice of $u$
midstream}. If you take $u = e^x$ in the first parts and then $u = \cos x$ (i.e.
$\dd v = e^x\,\dd x$) in the second, the second application simply undoes the
first: you will get $K = K$, the useless identity $0=0$, and you'll think the
method failed. It didn't -- you sabotaged it. The rule is iron: \emph{be
consistent}. If $u$ was the exponential the first time, $u$ is the exponential
every time. Consistency is what lets the integral genuinely cycle back rather
than cancel.
\end{trapbox}

\begin{happeningbox}
Why does the integral come back rather than terminating? Because the pair
$\{\sin x, \cos x\}$ is \emph{closed} under differentiation -- differentiating
either one stays inside the pair -- and $e^x$ is its own derivative. So no factor
ever simplifies to zero the way a polynomial eventually does; the process can only
\emph{rotate} through a finite set of shapes. After two rotations of the trig
factor you are back where you started, up to a sign and the boundary terms that
piled up along the way. The ``recurring integral'' is the visible footprint of
that rotation. Linear algebra has a name for this: $\sin x$ and $\cos x$ span a
two-dimensional space invariant under $\dv{}{x}$, and integration by parts is
tracing a loop around it.
\end{happeningbox}

\section{Tabular Integration: Parts on Autopilot}
\label{ch23:tabular}

When one factor is a polynomial that differentiates down to zero -- like $x^3$ or
$x^2$ -- and the other is something you can integrate over and over without effort
-- like $e^x$, $\sin x$, $\cos x$ -- repeated integration by parts becomes purely
mechanical. There is a bookkeeping shortcut, often called the \emph{tabular
method} or the \emph{DI method}, that does all the rounds at once and never makes
a sign error.

Here is the recipe. Make two columns. In the $D$ column (for
\emph{D}ifferentiate) put your polynomial $u$ and keep differentiating down to
$0$. In the $I$ column (for \emph{I}ntegrate) put the $\dd v$ function and keep
integrating, one row for each row of the $D$ column. Then draw diagonal arrows
from each entry in $D$ to the entry one row \emph{down} in $I$, attach alternating
signs $+,-,+,-,\dots$ starting with $+$, and read off the products.

\begin{example}[Tabular: $\int x^3 e^{x}\,\dd x$]
\label{ch23:ex-tabx3}
Differentiate $x^3$ down to zero; integrate $e^x$ (it never changes):
\begin{center}
\begin{tabular}{c c c c}
\toprule
sign & $D$ (differentiate) & & $I$ (integrate) \\
\midrule
$+$ & $x^3$ & $\searrow$ & $e^x$ \\
$-$ & $3x^2$ & $\searrow$ & $e^x$ \\
$+$ & $6x$ & $\searrow$ & $e^x$ \\
$-$ & $6$ & $\searrow$ & $e^x$ \\
& $0$ & & $e^x$ \\
\bottomrule
\end{tabular}
\end{center}
Multiply each $D$ entry by the $I$ entry one row down, with the sign on its left:
\[
\int x^3 e^x\,\dd x
= +x^3 e^x - 3x^2 e^x + 6x e^x - 6 e^x + C
= (x^3 - 3x^2 + 6x - 6)e^x + C.
\]
Compare this to the laborious three-rounds-of-parts we would otherwise do, and to
the reduction formula's $I_3$ in Example~\ref{ch23:ex-redxn} -- same answer, a
fraction of the writing.
\end{example}

\begin{whybox}
Tabular integration is not a separate technique -- it is ordinary parts with the
bookkeeping automated. Watch the columns line up with the formula. The first
diagonal $+x^3\cdot e^x$ is the boundary term $uv$. The next row down is exactly
the $-\int v\,\dd u$ that parts hands you, but instead of writing the integral
sign you immediately apply parts to \emph{it}, producing the next $uv$ term and
the next leftover -- and so on. The alternating signs are the chain of minus signs
from each application of the formula. The table terminates precisely because
differentiating the polynomial eventually gives $0$, killing the final leftover
integral. So the table is a flip-book of the recursion in
Example~\ref{ch23:ex-redxn}, all frames drawn at once.
\end{whybox}

\begin{example}[Tabular with a trig factor: $\int x^2\cos x\,\dd x$]
\label{ch23:ex-tabx2cos}
Differentiate $x^2$ to zero; integrate $\cos x \to \sin x \to -\cos x \to
-\sin x$:
\begin{center}
\begin{tabular}{c c c c}
\toprule
sign & $D$ & & $I$ \\
\midrule
$+$ & $x^2$ & $\searrow$ & $\cos x$ \\
$-$ & $2x$ & $\searrow$ & $\sin x$ \\
$+$ & $2$ & $\searrow$ & $-\cos x$ \\
& $0$ & & $-\sin x$ \\
\bottomrule
\end{tabular}
\end{center}
Read the diagonals with signs:
\[
\int x^2\cos x\,\dd x
= +x^2\sin x - 2x(-\cos x) + 2(-\sin x) + C
= x^2\sin x + 2x\cos x - 2\sin x + C.
\]
The signs in the answer come from \emph{two} sources -- the alternating $\pm$ of
the table and the minus signs that appear inside the $I$ column as $\cos$
integrates to $\sin$ to $-\cos$. Keep them separate and the table never lies.
\end{example}

\begin{trapbox}
Tabular integration in its simplest form requires the $D$ column to terminate at
$0$ -- that is, $u$ must be a polynomial. If \emph{neither} factor differentiates
to zero (as with $\int e^x\cos x\,\dd x$, where both columns cycle forever), the
plain table never ends. There is an extended tabular method for those cases: you
stop when a row reproduces the original integrand and then solve algebraically --
which is just the recurring-integral trick of Section~\ref{ch23:recurring} drawn
in a table. But for that case the by-hand ``solve for $K$'' method is usually
clearer, so reach for the table only when a polynomial guarantees a stop.
\end{trapbox}

\section{Putting It All Together: A Harder Example}
\label{ch23:harder}

One last integral, and this one needs taste, not just a recipe.

\begin{example}[$\int x\arctan x\,\dd x$]
\label{ch23:ex-xarctan}
Two factors: algebraic $x$ and inverse-trig $\arctan x$. LIATE puts I above A, so
$u = \arctan x$ (differentiate the nasty one) and $\dd v = x\,\dd x$. Then
\[
\dd u = \frac{1}{1+x^2}\,\dd x, \qquad v = \frac{x^2}{2}.
\]
Parts:
\[
\int x\arctan x\,\dd x
= \frac{x^2}{2}\arctan x - \frac12\int \frac{x^2}{1+x^2}\,\dd x.
\]
The leftover is a rational function. Do the polynomial-division trick:
$\dfrac{x^2}{1+x^2} = 1 - \dfrac{1}{1+x^2}$. Hence
\[
\int \frac{x^2}{1+x^2}\,\dd x = \int 1\,\dd x - \int\frac{1}{1+x^2}\,\dd x
= x - \arctan x.
\]
Assemble the final answer:
\[
\int x\arctan x\,\dd x
= \frac{x^2}{2}\arctan x - \frac12\big(x - \arctan x\big) + C
= \frac{x^2+1}{2}\arctan x - \frac{x}{2} + C.
\]
Three techniques in one problem: parts to break it open, polynomial division to
simplify the rational leftover, and the known antiderivative of $1/(1+x^2)$. This
is what real integration looks like -- a sequence of recognitions, each one
turning the problem into a slightly easier problem until nothing is left.
\end{example}

\begin{aestheticbox}
Step back and admire what integration by parts \emph{is}. Differentiation is a
one-way street: given a formula, you can always grind out its derivative.
Integration is the reverse, and the reverse of a mechanical process is generally
an art -- there is no algorithm that integrates every elementary function in
elementary terms (some, like $\int e^{-x^2}\,\dd x$, have no elementary
antiderivative at all). Integration by parts is one of the few genuinely
\emph{structural} tools we have: it does not compute an answer so much as
\emph{re-describe the problem}, trading one integral for another and trusting that
the trade leads somewhere better. Substitution undoes composition; parts undoes
multiplication. With those two reversals, plus a handful of algebraic identities,
you can integrate a startling fraction of the functions that arise in physics,
probability, and engineering. The rest -- the genuinely non-elementary ones --
are where the next chapters, and ultimately infinite series, take over.
\end{aestheticbox}

\begin{applicationbox}
Integration by parts is not a calculus-class curiosity; it is structural
machinery throughout science. In \textbf{probability}, the expected value of a
nonnegative random variable can be rewritten as
$E[X] = \int_0^\infty (1-F(x))\,\dd x$ by integrating $\int_0^\infty x f(x)\,\dd x$
by parts -- the boundary term and the survival function $1-F$ appear precisely
because of the $uv$ term. In \textbf{physics}, the entire framework of Lagrangian
mechanics and field theory rests on ``integrating by parts to move a derivative
off the variation'': the Euler--Lagrange equations come from doing exactly that to
$\int \frac{\partial L}{\partial \dot q}\,\delta\dot q\,\dd t$. In \textbf{signal
processing}, every integration-by-parts shifts a derivative from one factor onto
another, which is the seed of the theory of \emph{distributions} (the
``derivative'' of a non-differentiable signal is \emph{defined} by where parts
would have moved it). In \textbf{numerical analysis}, repeated parts produces
asymptotic expansions for integrals that are otherwise intractable. Learn the move
well; you will meet it for the rest of your mathematical life.
\end{applicationbox}

\section{Why It Works, One More Time}
\label{ch23:whyworks}

Here is the core idea plainly, stripped of all examples. Once you hold it you
will never be confused about when to reach for parts.

\begin{intuitionbox}
\textbf{The one-sentence summary.} Integration by parts moves a derivative from
one factor of a product to the other, paying a boundary term $uv$ for the move.
\end{intuitionbox}

To see that, rewrite the formula with the differentials spelled out. Starting
from $\int u\,v'\,\dd x = uv - \int u'\,v\,\dd x$, notice that on the left the
prime (the derivative) sits on $v$, and on the right it sits on $u$. The two
integrals
\[
\int u\,v'\,\dd x \quad\longleftrightarrow\quad \int u'\,v\,\dd x
\]
are the same product with the derivative \emph{relocated} from one factor to the
other, and the cost of relocation is the rectangle $uv$ evaluated at the
endpoints. That is the whole mechanism. You reach for parts whenever moving the
derivative makes the integrand simpler:
\begin{itemize}
\item Move the derivative \emph{onto} a polynomial repeatedly and it dies
($x^n \to \dots \to 0$); the integral terminates. (This is tabular integration.)
\item Move the derivative \emph{off} a log or inverse-trig function and its ugly
``no antiderivative'' problem dissolves into a simple rational factor. (This is
the $\int\ln x\,\dd x$ trick.)
\item Move the derivative around a closed family like $\{e^x,\sin x,\cos x\}$ and
it cycles back, letting you solve for the integral algebraically. (This is the
recurring trick.)
\end{itemize}
Every example in this chapter is one of these three motions. The formula is one
line; the judgment is in seeing which direction to push the derivative so the
problem gets smaller.

\begin{formalbox}
\textbf{Statement and proof, for the record.} Let $u,v:[a,b]\to\R$ be continuous
with continuous derivatives $u',v'$ on $[a,b]$. Then $uv$ is differentiable and,
by the product rule, $(uv)' = u'v + uv'$, which is continuous. By the Fundamental
Theorem of Calculus,
\[
\eval{u(x)v(x)}{a}^{b} = \int_a^b (uv)'\,\dd x = \int_a^b u'v\,\dd x + \int_a^b uv'\,\dd x.
\]
Rearranging gives the definite form
$\int_a^b uv'\,\dd x = \eval{uv}{a}^{b} - \int_a^b u'v\,\dd x$; dropping the
limits and writing $\dd v = v'\,\dd x$, $\dd u = u'\,\dd x$ gives the indefinite
form. The only hypothesis is that $u'$ and $v'$ are continuous so that the
Fundamental Theorem applies -- everything else is algebra. $\qed$
\end{formalbox}

\section{What integration by parts unlocked, limitations, and the road ahead}

\begin{connectionbox}
Integration by parts is the calculus equivalent of a balance scale: when an integral is too heavy
to lift directly, we transfer weight from one pan to the other until the problem balances.

Before Taylor's product-rule inversion, functions from different mathematical kingdoms could not
be integrated together. A polynomial multiplied by an exponential ($x^n e^x$), a power of $x$
multiplied by a trigonometric wave ($x^n \sin x$), or solitary transcendental functions lacking
algebraic partners ($\ln x$, $\arctan x$, $\arcsin x$) were dead ends for substitution. By setting
$\dd v = \dd x$, integration by parts manufactured an artificial partner from thin air, liberating
inverse functions and logarithms from computational limbo. Furthermore, by generating recursive
reduction formulas (such as the Wallis integrals for $\pi$), parts established a bridge between
continuous integration and discrete sequence analysis.

Yet parts works only if transferring the derivative reduces the algebraic complexity of the integrand.
When it does not, parts runs into two distinct failure modes:
\begin{enumerate}
  \item \emph{Infinite cycling and spiraling complexity:} While combinations like $e^x\sin x$ cycle
  cleanly and can be solved by algebraic rearrangement, other simple-looking products spiral into
  worse expressions. Consider $\int \frac{e^x}{x}\,\dd x$: setting $u = 1/x$ produces
  $-\int \frac{e^x}{x^2}\,\dd x$, increasing the negative power without mercy; setting $u = e^x$
  produces $\int e^x \ln x\,\dd x$. Parts cannot simplify functions where both factors resist
  polynomial annihilation.
  \item \emph{Rational quotients:} Integrands consisting of polynomial fractions, such as
  $\int \frac{x^3 + 5x - 2}{x^4 - 16}\,\dd x$ or $\int \frac{1}{x^2+1}\,\dd x$, completely defy
  integration by parts. Choosing $u$ as the numerator and $\dd v$ as the denominator (or vice versa)
  leaves higher-degree rational functions or complicated logarithms that refuse to simplify.
\end{enumerate}

To conquer rational quotients and intricate trigonometric combinations, Chapter~24 introduces
\emph{Partial Fractions and Trigonometric Integrals}. By deploying the Fundamental Theorem of Algebra,
Chapter~24 decomposes any rational function $P(x)/Q(x)$ into simple linear and quadratic fractions
whose antiderivatives are elementary logarithms and arctangents, providing a complete, deterministic
algorithm for the entire universe of rational functions.
\end{connectionbox}

\section*{Exercises}
\addcontentsline{toc}{section}{Exercises}

\begin{exercise}[Warm-ups]
\label{ch23:exr-warmup}
Evaluate each indefinite integral using integration by parts. State your choice of
$u$ and $\dd v$ each time.
\begin{enumerate}[label=(\alph*)]
\item $\displaystyle\int x\cos x\,\dd x$
\item $\displaystyle\int x e^{-x}\,\dd x$
\item $\displaystyle\int x\,\ln x\,\dd x$
\item $\displaystyle\int \arcsin x\,\dd x$
\end{enumerate}
\end{exercise}

\begin{exercise}[Two rounds and tabular]
\label{ch23:exr-tworounds}
\begin{enumerate}[label=(\alph*)]
\item Evaluate $\displaystyle\int x^2 e^{x}\,\dd x$ by applying parts twice.
\item Now redo it with the tabular method and confirm you get the same answer.
\item Use the tabular method to evaluate $\displaystyle\int x^3\sin x\,\dd x$.
\end{enumerate}
\end{exercise}

\begin{exercise}[Definite integrals]
\label{ch23:exr-definite}
Evaluate, simplifying fully:
\begin{enumerate}[label=(\alph*)]
\item $\displaystyle\int_0^1 x e^{x}\,\dd x$
\item $\displaystyle\int_1^{e} x^2\ln x\,\dd x$
\item $\displaystyle\int_0^{\pi} x\sin x\,\dd x$
\end{enumerate}
\end{exercise}

\begin{exercise}[The recurring trick]
\label{ch23:exr-recurring}
\begin{enumerate}[label=(\alph*)]
\item Evaluate $\displaystyle\int e^{x}\sin x\,\dd x$ by applying parts twice and
solving for the integral. (Be consistent in your choice of $u$.)
\item Evaluate $\displaystyle\int e^{2x}\cos(3x)\,\dd x$ the same way.
\item Explain in one or two sentences why choosing $u$ inconsistently makes the
method collapse to $0=0$.
\end{enumerate}
\end{exercise}

\begin{exercise}[Reduction formulas]
\label{ch23:exr-reduction}
\begin{enumerate}[label=(\alph*)]
\item Derive the reduction formula
$\displaystyle\int (\ln x)^n\,\dd x = x(\ln x)^n - n\int(\ln x)^{n-1}\,\dd x$.
\item Use it to evaluate $\displaystyle\int (\ln x)^2\,\dd x$.
\item Derive the reduction formula
$\displaystyle\int\cos^n x\,\dd x = \frac{1}{n}\cos^{n-1}x\,\sin x
+ \frac{n-1}{n}\int\cos^{n-2}x\,\dd x$, mirroring the $\sin^n$ derivation in
Example~\ref{ch23:ex-redsin}.
\end{enumerate}
\end{exercise}

\begin{exercise}[Combining with substitution]
\label{ch23:exr-combine}
Each of these needs a substitution \emph{first}, then parts. Find them.
\begin{enumerate}[label=(\alph*)]
\item $\displaystyle\int \cos(\sqrt{x})\,\dd x$ \quad (try $w=\sqrt{x}$).
\item $\displaystyle\int x^3 e^{x^2}\,\dd x$ \quad (try $w=x^2$, then parts).
\item $\displaystyle\int \sin(\ln x)\,\dd x$ \quad (try $w=\ln x$, then the
recurring trick).
\end{enumerate}
\end{exercise}

\begin{exercise}[Conceptual]
\label{ch23:exr-concept}
\begin{enumerate}[label=(\alph*)]
\item A student integrates $\int x e^x\,\dd x$ by choosing $u=e^x$,
$\dd v = x\,\dd x$ and gets a \emph{worse} integral. Explain, in terms of the two
principles in Section~\ref{ch23:choosing}, exactly what went wrong.
\item Explain why $\int \sec x\,\dd x$ cannot be cracked by the
``$u=\sec x,\ \dd v=\dd x$'' trick that worked for $\ln x$ and $\arctan x$. What is
different about $\sec x$? (You do not have to evaluate it.)
\item Why does the geometric ``two ways of slicing a rectangle'' proof in
Section~\ref{ch23:picture} require $u$ and $v$ to be increasing, while the
algebraic derivation does not?
\end{enumerate}
\end{exercise}

\begin{exercise}[A proof]
\label{ch23:exr-proof}
Prove the \textbf{Wallis recursion} for the definite integrals
$\displaystyle W_n = \int_0^{\pi/2}\sin^n x\,\dd x$: show that
$W_n = \dfrac{n-1}{n}\,W_{n-2}$ for $n\ge 2$. Begin from the reduction formula in
Example~\ref{ch23:ex-redsin} and explain carefully why the boundary term
$-\frac1n\sin^{n-1}x\cos x$ vanishes at \emph{both} endpoints $x=0$ and
$x=\pi/2$. Then compute $W_0$ and $W_1$, and use the recursion to find
$W_4$ and $W_5$.
\end{exercise}

\begin{exercise}[Going further]
\label{ch23:exr-further}
Use integration by parts to establish the formula for ``raising'' a power of
$\sec$:
\[
\int \sec^n x\,\dd x = \frac{\sec^{n-2}x\,\tan x}{n-1}
+ \frac{n-2}{n-1}\int \sec^{n-2}x\,\dd x \qquad (n\ge 2).
\]
\textit{Hint:} write $\sec^n x = \sec^{n-2}x\cdot\sec^2 x$, let
$\dd v = \sec^2 x\,\dd x$ (so $v=\tan x$), and use $\tan^2 x = \sec^2 x - 1$. You
will need to solve for the recurring integral, just as in
Section~\ref{ch23:recurring}.
\end{exercise}
